\documentclass[11pt,a4paper]{article}
% Common preamble for the hint sheets (\input from each file)
\usepackage[utf8]{inputenc}
\usepackage[T1]{fontenc}
\usepackage{lmodern}
\usepackage{amsmath,amssymb,amsthm,mathtools}
\usepackage[a4paper,margin=2.3cm]{geometry}
\usepackage{enumitem}
\usepackage{xcolor}
\usepackage[hidelinks]{hyperref}
\renewcommand{\leq}{\leqslant}
\renewcommand{\le}{\leqslant}
\renewcommand{\geq}{\geqslant}
\renewcommand{\ge}{\geqslant}
\setlength{\parindent}{0pt}
\setlength{\parskip}{0.35em}
% \hintstitle{Sheet title}{Class work / Homework}
\newcommand{\hintstitle}[2]{\begin{center}{\LARGE\bfseries Hints: #1\par}\vspace{0.5em}{\large #2\par}\end{center}
\noindent\textit{The hints to each problem are cumulative. Read one, try again, and open the next only if you are still stuck.}\par\medskip\hrule\medskip}
\newcommand{\problem}[1]{\par\bigskip{\large\bfseries Problem #1}\par\nopagebreak\smallskip}
\newcommand{\subpart}[1]{\par\smallskip\textbf{(#1)}\par\nopagebreak}
\newcommand{\hint}[1]{\par\smallskip\noindent\textcolor{black!70}{\textbf{Hint #1.}}\enspace}
\newcommand{\note}[1]{\par\smallskip\noindent\textit{Note.} #1\par}
\providecommand{\lt}{<}
\providecommand{\gt}{>}
\providecommand{\tightlist}{\setlength{\itemsep}{0pt}\setlength{\parskip}{0pt}}
\emergencystretch=3em

\begin{document}
\hintstitle{Methods of proof -- 6. Contradiction and proofs of impossibility}{Homework}

\problem{1}
\hint{1} Suppose $2^a,2^b,2^c$ with $a<b<c$ form an arithmetic progression: $2^a+2^c=2^{b+1}$. Divide by $2^a$ and look at parities.
\hint{2} $1+2^{c-a}$ is odd and greater than $1$, while $2^{b-a+1}$ is even. Contradiction (this is a $2$-adic valuation argument).

\problem{2}
\hint{1} Pairwise touching cubes pairwise intersect. Axis-parallel boxes that pairwise intersect have a common point: project to each coordinate axis and use Helly's theorem on the line (pairwise intersecting intervals have a common point).
\hint{2} So all cubes contain a common point $P$. For each cube $c+[0,1]^n\ni P$ and each coordinate $k$, at least one of $[P_k,P_k+\varepsilon]$, $[P_k-\varepsilon,P_k]$ lies in $[c_k,c_k+1]$ (for small $\varepsilon$). Hence each cube contains one of the $2^n$ small ``orthant boxes'' at $P$.
\hint{3} These orthant boxes have non-empty interior, and the cubes have disjoint interiors, so different cubes contain different orthant boxes: at most $2^n$ cubes. The $2^n$ cubes with a vertex at the origin show sharpness.

\problem{3}
\hint{1} You must show: for every $f$ there is a convergent sequence of pairwise distinct $a_i$ with $(f(a_i))$ convergent. Look for an infinite set on which both $x$ and $f(x)$ are bounded.
\hint{2} $\mathbb R=\bigcup_{m\ge1}\{x\in[-m,m]:\ |f(x)|\le m\}$ is a countable union, and $\mathbb R$ is uncountable, so one of these sets is infinite.
\hint{3} Take pairwise distinct points in it; by Bolzano--Weierstrass extract a convergent subsequence, and then a further subsequence along which the bounded values $f(a_i)$ converge.

\problem{4}
\hint{1} Encode a figure $F$ by the function $\mathbf 1_F(x,y)$ and use the rank: a rectangle $I\times J$ gives $\mathbf 1_I(x)\mathbf 1_J(y)$, a function of rank one. If $F$ is cut into $m$ rectangles, then $\mathbf 1_F$ (up to boundaries, e.g.\ as an integral operator on $L^2$) has rank at most $m$.
\hint{2} With the squares $I_i\times J_i$, $\mathbf 1_F=\sum_{i=1}^n\mathbf 1_{I_i}\otimes\mathbf 1_{J_i}$, and this has rank exactly $n$ provided the functions $\mathbf 1_{I_1},\dots,\mathbf 1_{I_n}$ are linearly independent, and likewise the $\mathbf 1_{J_i}$. So it suffices to show that intervals of lengths $1,\frac{11}{10},\dots,(\frac{11}{10})^{n-1}$ have independent indicators.
\hint{3} If $\sum c_i\mathbf 1_{I_i}=0$ with some $c_i\ne0$, view the intervals with $c_i\ne0$ as edges between their endpoints; the relation says the $c_i$ form a circulation, so these edges contain a cycle. Going around the cycle, the signed lengths add up to $0$: $\sum\pm(\frac{11}{10})^{k}=0$ over distinct exponents $k$.
\hint{4} This is impossible by an $11$-adic valuation: the terms have pairwise distinct $11$-adic valuations $k$, so the term with the smallest $k$ cannot be cancelled (multiply by $10^{n-1}11^{-k_{\min}}$ and reduce modulo $11$). Hence the rank is $n$ and at least $n$ rectangles are needed.

\problem{5}
\hint{1} Let $D=b^2-4ac$ (not a square, so $a\ne0$ and the form vanishes only at $(0,0)$). Take an odd prime $p\nmid aD$ with $\bigl(\frac Dp\bigr)=-1$. Show that if $p\mid Q(x,y)$ then $p\mid x$ and $p\mid y$; use $4aQ=(2ax+by)^2-Dy^2$.
\hint{2} Then $Q(x,y)=p^2Q(x/p,y/p)$, and by descent $v_p(Q(x,y))$ is even for $(x,y)\ne(0,0)$; hence $v_p(Q(x,y)^z)=z\,v_p(Q(x,y))$ is even too. So a number $N$ with $v_p(N)=1$ is not representable.
\hint{3} There are infinitely many primes with $\bigl(\frac Dp\bigr)=-1$ (quadratic reciprocity plus Dirichlet's theorem, or a direct argument). Choose $n$ of them, $p_1,\dots,p_n$, not dividing $2aD$, and by the Chinese remainder theorem find $N$ with $N+i\equiv p_i\pmod{p_i^2}$ for $i=1,\dots,n$. Then $v_{p_i}(N+i)=1$ and none of $N+1,\dots,N+n$ is representable.

\problem{6}
\hint{1} Look at the area $s(t)$ of the horizontal cross-section at height $t$. For the parallelepiped it is constant on $(0,H)$. What is it for a tetrahedron with a horizontal face?
\hint{2} If the tetrahedron has a horizontal face of area $A$ at height $h_0$ and apex at height $h_1$, its section at height $t$ between them is a similar triangle of area $A\bigl(\frac{t-h_1}{h_0-h_1}\bigr)^2$: a quadratic in $t$ with positive leading coefficient, and $0$ outside $[h_0,h_1]$.
\hint{3} On an interval between consecutive heights of vertices, the constant $s(t)$ would equal a sum of such quadratics, at least one of which is present (the section is non-empty). The coefficient of $t^2$ of the sum is positive: contradiction. (Equivalently, $s''=\sum s_T''>0$ there.)

\problem{7}
\hint{1} The answer is yes. Look among commuting nilpotent matrices of rank one, e.g.\ $A=E_{12}$ and $B=E_{13}$ (matrix units); check $AB=BA=0$.
\hint{2} Both have $A^2=B^2=0$ and rank $1$, so in a Jordan basis $f_1,f_2,f_3$ each is one $2\times2$ block on two consecutive basis vectors plus a zero block: $A$ sends one basis vector $f_{j+1}$ to $f_j$ and kills the other two, and $\operatorname{im}A=\langle f_j\rangle$.
\hint{3} Here $\operatorname{im}A=\operatorname{im}B=\langle e_1\rangle$, so in a common Jordan basis the blocks of $A$ and $B$ sit at the same place, and then $A$ and $B$ have the same matrix, i.e.\ $A=B$. But $A\ne B$.

\problem{8}
\hint{1} The image of a continuous map from $[0,1]^2$ is compact and connected, so for an injection it is a non-degenerate segment $[a,b]$. Remove a point.
\hint{2} Choose $P$ with $f(P)\in(a,b)$. The square minus $P$ is connected, but its image $[a,b]\setminus\{f(P)\}$ is not: contradiction.
\hint{3} A continuous bijection $[0,1]\to[0,1]^2$ would be a homeomorphism (compact to Hausdorff); removing an interior point disconnects $[0,1]$ but not the square. A continuous surjection exists (Peano curve), since surjectivity does not see this. The invariant in both cases is connectedness of the complement of a point (cut points).

\problem{9}
\hint{1} $\sqrt[3]2$ is a root of $x^3-2$; $\cos20^\circ$ satisfies $4c^3-3c=\cos60^\circ=\frac12$, i.e.\ $x=2c$ is a root of $x^3-3x-1$. A cubic over $\mathbb Q$ is irreducible iff it has no rational root.
\hint{2} By the rational root test the candidates are $\pm1,\pm2$ and $\pm1$ respectively, and none works; so both numbers have degree $3$.
\hint{3} Theorem: if $\alpha$ is constructible from the points $0,1$, then $\mathbb Q(\alpha)$ lies at the top of a tower $\mathbb Q=K_0\subset K_1\subset\dots\subset K_m$ with $[K_{i+1}:K_i]=2$, so $[\mathbb Q(\alpha):\mathbb Q]$ is a power of $2$ by the tower law. Doubling the cube needs $\sqrt[3]2$, trisecting $60^\circ$ needs $\cos20^\circ$; degree $3$ is not a power of $2$.

\problem{10}
\note{The quoted argument is in fact a correct proof by contradiction once two implicit steps are justified, so ``the error'' to be found is the omission of these justifications.}
\hint{1} Check each step under the hypothesis $\sum1/n=S<\infty$. Is regrouping a convergent series allowed? Does a strict termwise inequality between two convergent series give a strict inequality of the sums?
\hint{2} Both are true: regrouping a convergent series does not change its sum, and $\sum(a_k-b_k)>0$ if all $a_k-b_k>0$. So the argument is valid; what is missing is exactly these justifications, which is where convergence is used (the same manipulations without the hypothesis would be meaningless).
\hint{3} A direct proof: $s_{2^m}\ge1+\frac m2$ by grouping $\frac1{2^{j-1}+1}+\dots+\frac1{2^j}\ge\frac12$, so the partial sums are unbounded. It also gives the growth rate, which is why it is the one to submit.

\end{document}
